% explainer-source-hash: a04a073084a385412e9d3c61132e4cf9f1f315d82c412627f62fe74e5c87d8c7
% explainer-source-commit: dc7a1bb6930b81bc5a0e8ee9ee020d240cf03e21
% explainer-generated: 2026-07-16
% Regenerate with /explain-all (see WIRING.md). Do not edit this stamp by hand:
% it is how the toolchain knows whether this document still matches the code.
\documentclass[11pt,a4paper]{article}
\input{preamble}

\title{\textbf{prompt-evolution-ga}\\[0.3em]
\large A Deep Dive: evolving natural-language LLM prompts\\with a genetic algorithm and a compile-and-run fitness function}
\author{Explainer for Salman Adnan's portfolio}
\date{}

\begin{document}
\maketitle
\thispagestyle{empty}

\begin{keyidea}{What this document is}
An exhaustive, beginner-safe walkthrough of the \file{prompt-evolution-ga} project: what it
does, how every module works, what the committed results actually say, and where the code
and its own documentation disagree. It assumes no prior knowledge of genetic algorithms,
large language models, or C++ tooling. Every technical term is defined at first use.

This is a working document for the owner's own understanding and interview defense. It is
deliberately blunt about the project's weak points, because a flaw you learn about here is
a flaw you can defend in a room; one an interviewer finds first is not.
\end{keyidea}

\tableofcontents
\newpage

%=============================================================================
\section{The project in one page}
%=============================================================================

Large language models that write code are unreasonably sensitive to the wording of the
instructions you give them. Telling a model ``you are an expert competitive programmer,
consider edge cases'' can produce measurably different code from ``just give me the
code''. Finding the wording that works is normally a human sitting there, guessing,
re-running, and guessing again.

This project asks: \emph{can a genetic algorithm do that guessing automatically, and can we
prove whether it worked?}

\begin{jargon}{Large language model (LLM)}
A program trained on enormous amounts of text that, given some text, predicts plausible
text to follow. Modern ones (here, Google's Gemini) are good enough at this that they can
write working computer programs from an English description of a problem. They are not
databases and not calculators: the same question can yield different answers, and small
changes in how you ask can change the answer a lot.
\end{jargon}

\begin{jargon}{Prompt, and system prompt}
A \textbf{prompt} is the text you send an LLM. Many LLM interfaces split this in two: a
\textbf{system prompt}, which sets standing instructions and persona (``you are an expert
competitive programmer, always think about edge cases''), and a \textbf{user message},
which carries the actual task (``here is the problem, solve it''). This project holds the
user message fixed and evolves \emph{only the system prompt}. That is the whole point:
the task never changes, so any change in code quality is attributable to the wording of
the standing instructions.
\end{jargon}

\begin{jargon}{Genetic algorithm (GA)}
An optimisation method that copies biological evolution. You keep a \textbf{population} of
candidate solutions. Each is scored by a \textbf{fitness function}. The fitter ones are
more likely to be chosen as \textbf{parents}. Parents are combined (\textbf{crossover})
and randomly tweaked (\textbf{mutation}) to make \textbf{offspring}, which form the next
\textbf{generation}. Repeat. Over many generations the population drifts toward
higher fitness. Nothing here understands the problem; it only knows which candidates
scored better.
\end{jargon}

\subsection{The three design choices that define the project}

\begin{enumerate}
  \item \textbf{The chromosome is English prose.} In a textbook GA, a candidate solution is
  a fixed-length vector of numbers or bits. Here a candidate is an entire
  natural-language instruction prompt. This choice drives everything else, and it breaks
  the usual GA guarantees (Section~\ref{sec:operators}).

  \item \textbf{Fitness is objective and brutal.} A candidate prompt is not judged by
  asking another model ``is this a good prompt?''. It is judged by: send the prompt plus a
  real competitive-programming problem to Gemini, take the C++ it writes, \emph{actually
  compile it} with \code{g++}, \emph{actually run} the resulting binary against the
  problem's real test cases, and score the fraction of tests that pass. There is no room
  to argue with a failed test.

  \item \textbf{The genetic operators are themselves an LLM.} You cannot meaningfully
  ``swap the second half'' of two English paragraphs and get English. So crossover and
  mutation are performed by a second, stronger model (Gemini 2.5 Pro) which is asked to
  merge, extend, trim or rephrase the prompt text.
\end{enumerate}

\begin{keyidea}{The one-sentence summary}
A genetic algorithm whose population is English prompts, whose fitness is
``does the C++ this prompt produced actually compile and pass the tests'', and whose
crossover and mutation operators are a second LLM.
\end{keyidea}

\subsection{Provenance and honesty note}

The repository contains a course paper (Computational Intelligence), co-authored with
Ghufran Alvi; the implementation is Salman Adnan's work. Throughout this document, where a
claim in the README or the paper conflicts with what the committed code and data actually
show, \textbf{the code and data win}, and the conflict is called out in a red box. There
are several, and two of them are substantial.

\newpage
%=============================================================================
\section{Architecture: the whole system at a glance}
%=============================================================================

Two external systems sit outside the project's control: Google Cloud Vertex AI (which
serves the Gemini models) and the local machine's C++ toolchain (\code{g++}). Everything
else is ten Python files with no framework and no database.

\begin{figure}[H]
\centering
\resizebox{\textwidth}{!}{
\begin{tikzpicture}

  % --- explicit grid, so nothing can collide ---
  \node[box, minimum width=32mm] (main)      at ( 0.0,  0.0) {\textbf{main.py}\\orchestrator\\pre-flight, resume, holdout};
  \node[box, minimum width=24mm] (config)    at (-5.2,  0.0) {\textbf{config.py}\\every constant};

  \node[box, minimum width=30mm] (ga)        at ( 0.0, -2.8) {\textbf{ga.py}\\the GA loop\\Individual, selection,\\breed, fitness};
  \node[box, minimum width=24mm] (seeds)     at (-5.2, -2.8) {\textbf{seeds.py}\\15 hand-written\\seed prompts};
  \node[box, minimum width=24mm] (dataset)   at ( 5.2, -2.8) {\textbf{dataset.py}\\load + split\\APPS problems};

  \node[box, minimum width=28mm] (llm)       at (-3.0, -6.0) {\textbf{llm.py}\\Gemini calls:\\code gen + operators};
  \node[box, minimum width=28mm] (evaluator) at ( 3.0, -6.0) {\textbf{evaluator.py}\\compile, run,\\score};
  \node[box, minimum width=26mm] (cost)      at (-7.6, -6.0) {\textbf{cost\_tracker.py}\\tokens + dollars\\per call};

  \node[extbox, minimum width=30mm] (vertex) at (-3.0, -9.2) {Google Cloud Vertex AI\\(Gemini 2.5 Flash / Pro)};
  \node[extbox, minimum width=30mm] (gpp)    at ( 3.0, -9.2) {local toolchain\\\code{g++ -O2 -std=c++23}};

  \node[store, minimum width=18mm] (costs)   at (-7.6, -9.2) {Costs/};
  \node[store, minimum width=18mm] (results) at ( 7.6, -6.0) {results/\\run\_*/};

  \draw[flow] (config)    -- (main);
  \draw[flow] (main)      -- (ga);
  \draw[flow] (seeds)     -- (ga);
  \draw[flow] (dataset)   -- (ga);
  \draw[flow] (ga)        -- (llm);
  \draw[flow] (ga)        -- (evaluator);
  \draw[flow] (llm)       -- (vertex);
  \draw[flow] (evaluator) -- (gpp);
  \draw[flow] (llm)       -- (cost);
  \draw[dflow] (cost)     -- (costs);
  \draw[dflow] (ga.south east) -- node[lbl, pos=0.62] {checkpoints,\\eval logs} (results.west);

\end{tikzpicture}}
\caption{Module map. Solid arrows are calls; dashed arrows are writes to disk. Dashed
boxes are systems outside the project: the model API and the C++ compiler. Every fitness
number the project reports depends on both of them.}
\end{figure}

\subsection{Directory structure}

\begin{figure}[H]
\centering
\begin{forest} dirtree
[prompt-evolution-ga
  [main.py, label=right:{\rmfamily\small entry point: pre-flight, GA, holdout eval}]
  [ga.py, label=right:{\rmfamily\small the GA core (495 lines, the heart)}]
  [llm.py, label=right:{\rmfamily\small Vertex AI client, code gen, operators}]
  [evaluator.py, label=right:{\rmfamily\small compile, run, score}]
  [dataset.py, label=right:{\rmfamily\small APPS loading and splitting}]
  [seeds.py, label=right:{\rmfamily\small 15 seed prompts (only 5 are used)}]
  [config.py, label=right:{\rmfamily\small all hyperparameters}]
  [cost\_tracker.py, label=right:{\rmfamily\small token and dollar accounting}]
  [run\_baseline\_eval.py, label=right:{\rmfamily\small evolved vs baseline on the test split}]
  [visualize\_test\_results.py, label=right:{\rmfamily\small rebuilds the 5 comparison figures}]
  [Dataset/APPS/
    [train/, label=right:{\rmfamily\small 100 problems used here}]
    [test/, label=right:{\rmfamily\small 200 problems used here}]
  ]
  [results/
    [run\_20260429\_210411/, label=right:{\rmfamily\small tournament run}]
    [run\_20260430\_014526/, label=right:{\rmfamily\small roulette run}]
    [run\_20260430\_014530/, label=right:{\rmfamily\small elitism run}]
    [comparison/, label=right:{\rmfamily\small the 4-way test-set comparison + figures}]
  ]
  [Costs/, label=right:{\rmfamily\small per-run cost reports}]
  [corrections.tex, label=right:{\rmfamily\small loose patch notes for the paper}]
]
\end{forest}
\caption{Repository layout. Of 1{,}428 tracked files, 1{,}201 are the APPS dataset and most
of the rest are run artifacts; the actual program is ten Python files totalling roughly
2{,}300 lines.}
\end{figure}

\newpage
%=============================================================================
\section{The evolutionary loop}
\label{sec:loop}
%=============================================================================

This is the centre of the project. Everything else feeds it.

\subsection{The cycle}

\begin{figure}[H]
\centering
\resizebox{0.86\textwidth}{!}{
\begin{tikzpicture}

  % --- left column: the main descent; right column: the breeding return path ---
  \node[box, minimum width=40mm] (seed)      at ( 0.0,   0.0) {\textbf{Seed population}\\5 hand-written prompts\\\code{seeds.py}};
  \node[box, minimum width=40mm] (evalpop)   at ( 0.0,  -2.7) {\textbf{Evaluate fitness}\\each prompt vs 65 problems\\\code{compute\_fitness}};
  \node[dec, minimum width=26mm] (converged) at ( 0.0,  -5.6) {converged?\\CV $<$ 0.01\\for 3 gens};
  \node[box, minimum width=40mm] (elite)     at ( 0.0,  -8.5) {\textbf{Carry elites}\\top \code{elitism} individuals\\pass through unchanged};
  \node[box, minimum width=40mm] (select)    at ( 0.0, -11.2) {\textbf{Selection}\\tournament or roulette\\pick parents $p_1$, $p_2$};

  \node[box, minimum width=40mm] (cross)     at ( 7.4, -11.2) {\textbf{Crossover} (rate 0.8)\\LLM merges $p_1$ and $p_2$\\else copy one parent};
  \node[box, minimum width=40mm] (mutate)    at ( 7.4,  -8.5) {\textbf{Mutation}\\inject 0.25 / delete 0.15\\/ rephrase 0.10 (LLM)};
  \node[box, minimum width=40mm] (offspring) at ( 7.4,  -5.6) {\textbf{Offspring}\\fill population to 5\\then evaluate fitness};

  \node[box, minimum width=36mm, fill=white, draw=warnred, text=warnred]
       (done) at (-6.6, -8.5) {\textbf{Return best}\\then 200-problem\\holdout evaluation};

  \draw[flow] (seed)    -- (evalpop);
  \draw[flow] (evalpop) -- (converged);
  \draw[flow] (converged) -- node[lbl, pos=0.5] {no} (elite);
  \draw[flow] (elite)   -- (select);
  \draw[flow] (select)  -- (cross);
  \draw[flow] (cross)   -- (mutate);
  \draw[flow] (mutate)  -- (offspring);

  % return path: straight up the right column, then left into evaluate. Crosses nothing.
  \draw[flow] (offspring.north) -- node[lbl, pos=0.72] {next\\generation} ($(offspring.north)+(0,1.4)$) -- (evalpop.east);

  % exit: out to the left, clear of every box in the centre column
  \draw[flow] (converged.west) -- node[lbl, pos=0.62] {yes, or\\gen $=$ 12} (-6.6,-5.6) -- (done.north);

\end{tikzpicture}}
\caption{The generational cycle, as implemented in \code{ga.run}. The population is fully
replaced every generation except for the elites. With the shipped settings the loop runs
12 generations over a population of 5.}
\end{figure}

\subsection{The loop in pseudocode}

\begin{pseudocode}{ga.run: the generational loop (condensed from ga.py, lines 312-495)}
\Step{\textbf{input:} seed\_prompts (5), problems (65), scheme, elitism}
\Step{\textbf{output:} best individual, fitness history}
\Step{}
\Step{rng $\leftarrow$ Random(RANDOM\_SEED = 42)}
\Step{select $\leftarrow$ roulette\_select if scheme = "roulette" else tournament\_select}
\Step{}
\Step{\textbf{if} resuming \textbf{then}}
\Indent{population $\leftarrow$ checkpoint.population;  restore rng state}
\Step{\textbf{else}}
\Indent{population $\leftarrow$ [Individual(p) for p in seed\_prompts]}
\Indent{\textbf{for each} ind in population:}
\Indent{\hspace*{2em} \textbf{if} a seed checkpoint exists: restore its fitness (skip re-evaluating)}
\Indent{\hspace*{2em} \textbf{else}: ind.fitness $\leftarrow$ compute\_fitness(ind, problems); write checkpoint}
\Step{}
\Step{stagnant $\leftarrow$ 0}
\Step{\textbf{for} gen = 1 \textbf{to} GENERATIONS (12):}
\Indent{sort population by fitness, descending}
\Indent{stats $\leftarrow$ \{best, mean, worst, std, cv\} over the population}
\Indent{append stats to history;  save checkpoint of generation gen-1}
\Indent{}
\Indent{\textbf{if} 0 $\le$ cv $<$ 0.01 \textbf{then}  // population has gone homogeneous}
\Indent{\hspace*{2em} stagnant $\leftarrow$ stagnant + 1}
\Indent{\hspace*{2em} \textbf{if} stagnant $\ge$ 3: \textbf{break}   // converged, stop early}
\Indent{\textbf{else} stagnant $\leftarrow$ 0}
\Indent{}
\Indent{new\_pop $\leftarrow$ population[0 : elitism]        // elites survive untouched}
\Indent{off\_idx $\leftarrow$ 0}
\Indent{\textbf{while} |new\_pop| $<$ POPULATION\_SIZE (5):}
\Indent{\hspace*{2em} $p_1 \leftarrow$ select(population, rng)}
\Indent{\hspace*{2em} $p_2 \leftarrow$ select(population, rng)}
\Indent{\hspace*{2em} allow\_mut $\leftarrow$ (off\_idx $\ge$ CROSSOVER\_ONLY\_SLOTS = 2)}
\Indent{\hspace*{2em} new\_pop.append( breed($p_1$, $p_2$, rng, allow\_mut) )}
\Indent{\hspace*{2em} off\_idx $\leftarrow$ off\_idx + 1}
\Indent{}
\Indent{\textbf{for each} child in new\_pop[elitism:]:      // elites keep their old score}
\Indent{\hspace*{2em} child.fitness $\leftarrow$ compute\_fitness(child, problems)}
\Indent{}
\Indent{population $\leftarrow$ new\_pop;  sort descending}
\Step{}
\Step{append final stats to history}
\Step{\textbf{return} population[0], history}
\end{pseudocode}

\begin{keyidea}{The two details in that loop worth remembering}
\textbf{Crossover-only slots.} The first two offspring of every generation
(\code{CROSSOVER\_ONLY\_SLOTS = 2}) are bred with mutation switched off. This deliberately
reserves part of each generation for pure recombination, so mutation cannot scramble every
child.

\textbf{Elites are not re-evaluated.} An elite carries its old fitness number forward
rather than being re-scored. Since the LLM is called at temperature 0 this is defensible,
but it does mean an elite's fitness is a measurement from an earlier generation, never
refreshed.
\end{keyidea}

\subsection{The Individual}

\begin{lstlisting}[language=Python, caption={ga.py, lines 49-58. The chromosome and its provenance.}]
@dataclass
class Individual:
    prompt: str                     # THE CHROMOSOME: an English instruction prompt
    fitness: float = -1.0           # mean pass rate; -1.0 means "not yet evaluated"
    uid: str = field(default_factory=lambda: uuid.uuid4().hex[:8])
    parent_a: str = ""              # uid of parent A  \
    parent_b: str = ""              # uid of parent B   |  genealogy
    operators: str = ""             # e.g. "crossover+inject"  |
    intermediate_steps: dict = field(default_factory=dict)  # /  prompt after each operator
    failure_examples: list = field(default_factory=list)    # feeds the inject mutation
\end{lstlisting}

\begin{jargon}{Dataclass}
A Python shorthand that turns a list of named fields into a class, generating the
boilerplate constructor automatically. \code{@dataclass} above the class is what triggers
it. \code{field(default\_factory=...)} is needed for defaults that must be freshly created
per instance (a new random \code{uid}, a new empty list), rather than shared between all
instances.
\end{jargon}

The last four fields are the project's most interesting engineering idea. Because
LLM-mediated operators do not preserve structure, an offspring can silently wander away
from what its parents encoded. \code{parent\_a}, \code{parent\_b}, \code{operators} and
\code{intermediate\_steps} record the full derivation of every individual, so you can
reconstruct after the fact whether an operator helped or just scrambled a good prompt.
\code{intermediate\_steps} stores the prompt text after \emph{each} operator in the chain
(\code{after\_crossover}, \code{after\_inject}, \code{after\_delete},
\code{after\_rephrase}), so a bad step can be pinpointed. Section~\ref{sec:genealogy}
covers what happened when this data was actually used.

\newpage
%=============================================================================
\section{Selection: three schemes}
%=============================================================================

Selection decides which individuals get to be parents. The project implements two
mechanisms and exposes three schemes over them via \code{--scheme}.

\begin{figure}[H]
\centering
\resizebox{\textwidth}{!}{
\begin{tikzpicture}

  % tournament
  \node[font=\bfseries\small, text=inkblue] at (0,2.6) {Tournament selection (TOURNAMENT\_SIZE = 3)};
  \node[box, minimum width=8mm] (a1) at (-2.6,1.6) {0.62};
  \node[box, minimum width=8mm] (a2) at (-1.4,1.6) {0.71};
  \node[box, minimum width=8mm] (a3) at (-0.2,1.6) {0.55};
  \node[box, minimum width=8mm] (a4) at (1.0,1.6) {0.68};
  \node[box, minimum width=8mm] (a5) at (2.2,1.6) {0.60};
  \node[lbl] at (4.4,1.6) {population of 5};

  \node[box, minimum width=8mm, draw=accent, very thick] (b1) at (-1.4,0.2) {0.71};
  \node[box, minimum width=8mm, draw=accent, very thick] (b2) at (-0.2,0.2) {0.55};
  \node[box, minimum width=8mm, draw=accent, very thick] (b3) at (1.0,0.2) {0.68};
  \node[lbl] at (4.4,0.2) {sample 3 at random};

  \node[box, minimum width=8mm, fill=softgrey, draw=okgreen, very thick] (w) at (-1.4,-1.2) {0.71};
  \node[lbl] at (4.4,-1.2) {highest fitness wins};

  \draw[flow] (a2) -- (b1); \draw[flow] (a3) -- (b2); \draw[flow] (a4) -- (b3);
  \draw[flow] (b1) -- (w);

  % roulette
  \node[font=\bfseries\small, text=inkblue] at (11,2.6) {Roulette selection (fitness-proportional)};

  \fill[accent!70] (8.2,0.2) rectangle (10.0,0.9);
  \fill[accent!45] (10.0,0.2) rectangle (12.1,0.9);
  \fill[accent!25] (12.1,0.2) rectangle (13.7,0.9);
  \fill[accent!12] (13.7,0.2) rectangle (14.9,0.9);
  \draw[draw=inkblue] (8.2,0.2) rectangle (14.9,0.9);

  \node[font=\scriptsize, text=inkblue] at (9.1,0.55) {ind A};
  \node[font=\scriptsize, text=inkblue] at (11.05,0.55) {ind B};
  \node[font=\scriptsize, text=inkblue] at (12.9,0.55) {ind C};
  \node[font=\scriptsize, text=inkblue] at (14.3,0.55) {D};

  \node[lbl, align=center] at (11.5,1.5) {each individual's slice is proportional to its fitness};

  \draw[flow, draw=warnred, very thick] (11.6,-0.9) -- (11.6,0.15);
  \node[lbl, text=warnred] at (11.6,-1.3) {$r \sim \mathrm{Uniform}(0,\ \sum \mathrm{fitness})$ lands here $\rightarrow$ ind B};

\end{tikzpicture}}
\caption{The two selection mechanisms. Tournament applies constant pressure regardless of
the fitness \emph{scale}; roulette's pressure collapses when everyone scores similarly,
which is exactly what happens here (all fitness values sit near 0.65).}
\end{figure}

\subsection{The three schemes, and what actually distinguishes them}

\begin{lstlisting}[language=Python, caption={ga.py lines 329-333, plus main.py line 357. How --scheme maps onto behaviour.}]
# in ga.run:
scheme  = selection_scheme.lower()
select  = roulette_select if scheme == "roulette" else tournament_select
elitism = elitism_count if elitism_count is not None else config.ELITISM_COUNT

# in main.py, the only place elitism_count is ever passed:
elitism_count = 1 if args.scheme == "elitism" else None
\end{lstlisting}

So the three schemes decompose into just two independent switches:

\begin{table}[H]
\centering
\begin{tabular}{llll}
\toprule
\code{--scheme} & selection function & elitism & note \\
\midrule
\code{tournament} (default) & \code{tournament\_select} & \code{config.ELITISM\_COUNT} = 0 & \\
\code{roulette}            & \code{roulette\_select}   & \code{config.ELITISM\_COUNT} = 0 & \\
\code{elitism}             & \code{tournament\_select} & 1 (hardcoded in \code{main.py}) & tournament + elitism \\
\bottomrule
\end{tabular}
\caption{The \code{elitism} scheme is not a third selection mechanism. It is the
\code{tournament} scheme with one elite. The paper's own corrections file confirms this
reading.}
\end{table}

\begin{honest}{The scheme names oversell the experiment}
Presented as ``three selection schemes'', the experiment is really a $2 \times 1$ design:
two selection mechanisms (tournament, roulette) at elitism 0, plus one extra cell
(tournament at elitism 1). Roulette-with-elitism was never run. So the results cannot
separate ``elitism helps'' from ``elitism helps \emph{tournament}'', and there is no way to
tell whether the elitism run's distinct behaviour comes from the elite or from an
interaction. With a population of 5 and one run per cell, no scheme comparison here can
carry statistical weight; there is no repetition to average over.
\end{honest}

\subsection{Roulette selection, and its guard}

\begin{lstlisting}[language=Python, caption={ga.py, lines 245-256.}]
def roulette_select(population, rng):
    fits  = [max(ind.fitness, 0.0) for ind in population]   # clamp: -1.0 sentinel -> 0.0
    total = sum(fits)
    if total == 0:
        return rng.choice(population)                       # degenerate: pick uniformly
    r   = rng.uniform(0, total)
    cum = 0.0
    for ind, f in zip(population, fits):
        cum += f
        if cum >= r:
            return ind
    return population[-1]                                   # float-rounding safety net
\end{lstlisting}

Two defensive details worth noticing. The \code{max(ind.fitness, 0.0)} clamp matters
because \code{-1.0} is the project's sentinel for ``infrastructure failed, this is not a
real score''; without the clamp a negative would corrupt the running total and could make
the cumulative sum non-monotonic. The trailing \code{return population[-1]} catches the
case where floating-point rounding leaves \code{cum} a hair below \code{r} after the last
individual. Neither is exotic, but both are the kind of thing that is quietly wrong in most
first drafts.

\newpage
%=============================================================================
\section{The operators: crossover and mutation by LLM}
\label{sec:operators}
%=============================================================================

\begin{keyidea}{Why the operators had to be an LLM}
A classic GA chromosome is a fixed-length vector, so crossover is ``take the first $k$
genes from A and the rest from B'', and mutation is ``flip a random bit''. Both preserve
validity by construction: the result is always a well-formed chromosome.

An English paragraph has no such structure. Splicing the first half of one paragraph onto
the second half of another produces something that is not a valid instruction prompt. So
this project delegates the operators to Gemini 2.5 Pro and asks it, in English, to merge or
edit the prompts. The result is always fluent English. It is \emph{not} always related to
its parents, and that is the trade (see the honest box at the end of this section).
\end{keyidea}

\subsection{breed(): the offspring pipeline}

\begin{figure}[H]
\centering
\resizebox{\textwidth}{!}{
\begin{tikzpicture}

  \node[box, minimum width=22mm] (pa) at (0.0,  0.9) {parent A\\prompt};
  \node[box, minimum width=22mm] (pb) at (0.0, -0.9) {parent B\\prompt};

  \node[dec, minimum width=18mm] (d1) at (3.4, 0.0) {rng $<$ 0.8?};

  \node[box, minimum width=26mm] (xo)   at (7.4,  1.5) {\textbf{llm.crossover}\\Gemini 2.5 Pro\\merges both};
  \node[box, minimum width=26mm] (copy) at (7.4, -1.5) {\textbf{copy}\\pick one parent\\at random};

  \node[dec, minimum width=18mm] (d2) at (11.4, 0.0) {mutation\\allowed?};

  \node[box, minimum width=30mm] (inj) at (15.8,  3.0) {\textbf{inject} (p $=$ 0.25)\\add a strategy,\\conditioned on failures};
  \node[box, minimum width=30mm] (del) at (15.8,  0.0) {\textbf{delete} (p $=$ 0.15)\\drop the weakest\\sentence};
  \node[box, minimum width=30mm] (rep) at (15.8, -3.0) {\textbf{rephrase} (p $=$ 0.10)\\same meaning,\\new words};

  \node[box, minimum width=26mm, draw=okgreen] (child) at (20.6, 0.0) {\textbf{offspring}\\+ parents, ops,\\intermediate\_steps};

  \draw[flow] (pa) -- (d1); \draw[flow] (pb) -- (d1);
  \draw[flow] (d1) -- node[lbl, pos=0.45] {yes} (xo);
  \draw[flow] (d1) -- node[lbl, pos=0.45] {no} (copy);
  \draw[flow] (xo) -- (d2); \draw[flow] (copy) -- (d2);
  \draw[flow] (d2) -- node[lbl, pos=0.5] {yes} (inj);
  \draw[flow] (d2) -- (del);
  \draw[flow] (d2) -- (rep);
  \draw[flow] (inj) -- (child); \draw[flow] (del) -- (child); \draw[flow] (rep) -- (child);

  % bypass: routed well below the rephrase box so the label sits in clear space
  \draw[dflow] (d2.south) -- (11.4,-5.2)
      -- node[lbl, pos=0.5] {no (first 2 offspring):\\straight to offspring} (20.6,-5.2) -- (child.south);

\end{tikzpicture}}
\caption{\code{ga.breed}, lines 259-309. The three mutations are independent coin flips
applied in sequence, not a choice of one: an offspring can receive all three, or none.}
\end{figure}

The three mutation probabilities (0.25, 0.15, 0.10) are independent, so the chance an
offspring escapes mutation entirely is $(1-0.25)(1-0.15)(1-0.10) \approx 0.574$, and the
chance it gets all three is $0.25 \times 0.15 \times 0.10 = 0.00375$. The expected number of
mutations per mutable offspring is $0.25+0.15+0.10 = 0.5$.

\subsection{The empty-result guard, repeated four times}

\begin{lstlisting}[language=Python, caption={ga.py lines 271-298. The pattern that keeps a failed LLM call from destroying a chromosome.}]
if rng.random() < config.CROSSOVER_RATE:
    res    = llm.crossover(pa.prompt, pb.prompt)
    prompt = res if res.strip() else rng.choice([pa.prompt, pb.prompt])  # fall back
    ops.append("crossover")
else:
    prompt = rng.choice([pa.prompt, pb.prompt])
    ops.append("copy")

steps["after_crossover"] = prompt

if allow_mutation:
    if rng.random() < config.MUT_INJECT_PROB:
        res    = llm.mutate_inject(prompt, failure_examples=combined_fails)
        prompt = res if res.strip() else prompt      # empty -> keep previous prompt
        ops.append("inject")
        steps["after_inject"] = prompt
    # ... delete and rephrase repeat the identical shape
\end{lstlisting}

Every operator is wrapped in the same guard: if the model returns nothing usable, keep what
you had. This is the right instinct. A silent empty string from an API would otherwise
propagate an empty prompt into the population, and an empty prompt would score near zero
and pollute the generation.

\begin{honest}{The guard checks that output exists, not that it is complete}
\code{res.strip()} is truthy for \emph{any} non-blank text, including text the model was cut
off mid-word producing. That is not hypothetical here: it is exactly what happened, and it
is documented in Section~\ref{sec:truncation}. The right check is a completeness check
(does the prompt end at a sentence boundary? was the response truncated by the token cap?),
and the code does not do one.

Note also that \code{ops.append("inject")} runs \emph{unconditionally} once the coin flip
passes, even when the call came back empty and the prompt was left unchanged. So the
\code{operators} string can claim an individual was injected when nothing was injected. Any
later analysis of ``which operators helped'' inherits that error.
\end{honest}

\subsection{Failure-conditioned mutation: the best idea in the project}

The \code{inject} mutation is not blind. It is fed real failures from the parents' most
recent evaluation, so the model is asked to add a strategy targeting problems the parents
\emph{actually got wrong}.

\begin{lstlisting}[language=Python, caption={llm.py, lines 281-303. Building the failure-conditioned mutation prompt.}]
def mutate_inject(prompt: str, failure_examples: list | None = None) -> str:
    hint = ""
    if failure_examples:
        lines = []
        for f in failure_examples[:3]:
            line = f"  - Problem {f['problem_id']} (rating {f.get('rating','?')}): reason={f['reason']}"
            if f.get("compile_error"):
                line += f", compile_error={f['compile_error'][:80]!r}"
            elif f.get("expected"):
                line += f", expected={f['expected'][:60]!r}, actual={f['actual'][:60]!r}"
            lines.append(line)
        hint = ("\n\nRecent failure examples (address at least one):\n" + "\n".join(lines)
                + "\nAdd a strategy that targets these specific failure patterns.")

    user = ("Add one new, specific coding strategy to this instruction prompt. "
            "Examples: 'validate your output against the examples', "
            "'consider time complexity', 'check for off-by-one errors'."
            f"{hint}\n\nOutput only the modified prompt.\n\nPrompt:\n{prompt}")
\end{lstlisting}

\begin{keyidea}{Why this is a genuinely good design}
This turns a blind random mutation into a directed one. A standard GA mutation has no idea
\emph{why} an individual is weak; it perturbs at random and lets selection sort it out. Here
the mutation operator gets told ``your parents failed problem train/2015 with a compile
error, and problem test/0042 by printing 3 where 4 was expected''. That is a gradient
signal smuggled into a gradient-free method, and it is the closest thing the project has to
a genuinely novel idea.
\end{keyidea}

The failures are chosen by \code{\_pick\_failures} (\code{ga.py}, lines 217-228), which
sorts failures by a priority ranking (compile failures first, then total failures, then
partial passes) and keeps at most five. \code{breed} then merges both parents' failure
lists, de-duplicating by \code{problem\_id}, and caps at five.

\begin{honest}{\_pick\_failures has a logic bug that defeats its own diversity intent}
\begin{lstlisting}[language=Python]
for f in fails:
    if f["reason"] not in seen or len(picked) < 3:   # <-- the "or" is the bug
        picked.append(f)
        seen.add(f["reason"])
    if len(picked) >= 5:
        break
\end{lstlisting}
The evident intent is ``take one example per distinct failure reason, for variety''. But
\code{or len(picked) < 3} short-circuits that: while fewer than three are picked, the
condition is true \emph{regardless} of whether the reason was already seen. Since the list
is pre-sorted by reason priority, the first three picks are typically all the
\emph{same} reason (three compile failures, say). Only picks four and five apply the
distinct-reason rule. The operator therefore usually sees a less varied failure sample than
intended. Changing \code{or} to \code{and} would not fix it either; the loop needs
restructuring. This is invisible at runtime because nothing checks it, and it degrades the
project's best idea.
\end{honest}

\subsection{The output directives}

Both the code generator and the operators are wrapped in blunt, almost shouted formatting
directives (\code{\_CODE\_GEN\_DIRECTIVE} and \code{\_OPERATOR\_DIRECTIVE},
\code{llm.py} lines 45-69), forbidding markdown fences, preamble and commentary. This is
defensive prompt engineering: the pipeline needs raw C++ or raw prompt text, and any
``Here is your solution:'' preamble would be fed to the compiler.

Notice the layered defence: the directive tells the model not to emit code fences, and
\emph{then} \code{evaluator.extract\_code} strips code fences anyway. Belt and braces, and
correct: the directive reduces the failure rate, the extractor handles the residue.

\begin{honest}{The operator directives are a chromosome contaminant}
\code{\_CODE\_GEN\_DIRECTIVE} is appended to every candidate prompt at evaluation time, so
every individual is really ``evolved prompt + a fixed 700-character block of shouting''.
The block is constant across all individuals, so it cannot bias the \emph{comparison}
between prompts. But it does mean the fitness landscape being searched is the landscape
\emph{given} that directive: the evolved prompts are not standalone artifacts, and quoting
one in a paper without the directive misrepresents what was measured. The README does not
mention this, and the paper's Table of parameters does not either.
\end{honest}

\newpage
%=============================================================================
\section{Fitness: the compile-and-run pipeline}
%=============================================================================

This is the part of the project that makes it defensible. Fitness is not an opinion.

\begin{figure}[H]
\centering
\resizebox{\textwidth}{!}{
\begin{tikzpicture}

  % --- row 1, left to right: generate the code and compile it ---
  \node[box, minimum width=26mm]    (in)      at ( 0.0,  0.0) {candidate prompt\\(the chromosome)\\$+$ APPS problem};
  \node[box, minimum width=28mm]    (call)    at ( 4.2,  0.0) {\textbf{llm.generate\_code}\\system $=$ prompt $+$ directive\\user $=$ question $+$ rules};
  \node[extbox, minimum width=26mm] (gemini)  at ( 4.2,  2.9) {Gemini 2.5 Flash\\temperature 0.0\\thinking budget 512};
  \node[box, minimum width=24mm]    (extract) at ( 8.4,  0.0) {\textbf{extract\_code}\\strip code fences,\\else use raw text};
  \node[box, minimum width=24mm]    (compile) at (12.2,  0.0) {\textbf{\_compile}\\\code{g++ -O2 -std=c++23}\\15 s timeout};
  \node[dec, minimum width=18mm]    (ok)      at (15.8,  0.0) {compiled?};
  \node[box, minimum width=26mm, draw=warnred] (fail) at (19.8, 0.0) {score $=$ 0.0\\reason $=$ \code{compile\_fail}};

  % --- row 2, right to left: run the binary and score it ---
  \node[box, minimum width=28mm] (run)   at (15.8, -4.8) {\textbf{\_run\_one}, per test case\\pipe stdin, capture stdout\\2 s timeout, own process group};
  \node[box, minimum width=26mm] (cmp)   at (11.0, -4.8) {\textbf{\_match}\\strip trailing whitespace\\per line, compare exactly};
  \node[dec, minimum width=20mm] (early) at ( 6.4, -4.8) {3 consecutive\\failures?};
  \node[box, minimum width=28mm, draw=okgreen] (score) at ( 1.6, -4.8) {\textbf{score} $=$ passed / total\\reason $\in$ \{perfect,\\partial\_pass, all\_fail, ...\}};

  \draw[flow] (in)      -- (call);
  \draw[flow] (gemini)  -- (call);
  \draw[flow] (call)    -- (extract);
  \draw[flow] (extract) -- (compile);
  \draw[flow] (compile) -- (ok);
  \draw[flow] (ok) -- node[lbl, pos=0.45] {no} (fail);
  \draw[flow] (ok) -- node[lbl, pos=0.5] {yes} (run);
  \draw[flow] (run)   -- (cmp);
  \draw[flow] (cmp)   -- (early);
  \draw[flow] (early) -- node[lbl, pos=0.5] {yes:\\stop early} (score);

  % loop back for the next test case, routed clear underneath both boxes
  \draw[flow] (early.south) -- (6.4,-7.4)
      -- node[lbl, pos=0.5] {no: next test case} (15.8,-7.4) -- (run.south);

\end{tikzpicture}}
\caption{One fitness measurement for one prompt against one problem
(\code{evaluator.run\_tests\_verbose}). An individual's fitness is the mean of this over all
65 evolution problems, computed 5-way in parallel.}
\end{figure}

\subsection{Sandboxing: the details that matter}

\begin{lstlisting}[language=Python, caption={evaluator.py, lines 12-19 and 88-106. Killing a runaway program properly.}]
def _kill_group(proc):
    try:
        os.killpg(os.getpgid(proc.pid), signal.SIGKILL)   # kill the whole process GROUP
    except (ProcessLookupError, PermissionError):
        try: proc.kill()                                   # fall back to just the child
        except Exception: pass

# every subprocess is launched with:
proc = subprocess.Popen([...], preexec_fn=os.setsid, ...)  # new session => new process group
\end{lstlisting}

\begin{jargon}{Process group, and why \code{proc.kill()} is not enough}
When a program starts other programs, they normally form a family tree. \code{proc.kill()}
kills only the one process you launched. If the generated C++ had spawned children, or if
the compiler forks (which \code{g++} does: it is a driver that runs \code{cc1plus},
\code{as} and \code{ld}), killing the parent leaves orphans running and consuming CPU
forever. \code{os.setsid} puts each launched program into its own \emph{process group}, and
\code{os.killpg} kills every process in that group at once. Over three runs of roughly
13{,}000 model calls each producing a compile and up to a few dozen executions, that
distinction is the difference between a machine that survives the night and one that does
not.
\end{jargon}

Other real defences here: a 15 second compile timeout and a 2 second execution timeout
(\code{config.COMPILE\_TIMEOUT}, \code{config.EXEC\_TIMEOUT}); \code{stderr} redirected to
\code{DEVNULL} so a chatty program cannot fill a pipe buffer and deadlock; captured
input/output truncated to 500 characters (\code{\_MAX\_IO\_LEN}) before being logged, so a
program that prints a megabyte cannot bloat the JSONL logs; and the temporary source and
binary files cleaned up in a \code{finally} block.

\begin{honest}{This is a timeout, not a sandbox}
The generated C++ is compiled and executed on the host with the user's own privileges and
no isolation: no container, no seccomp filter, no user namespace, no filesystem or network
restriction. A model that emitted code calling \code{system("rm -rf ~")}, or opening a
socket, would be obeyed. The 2 second timeout bounds \emph{how long} arbitrary code runs,
not \emph{what} it can do.

In practice the risk is low: the inputs are benign competitive-programming problems and the
model has no motive. But ``we ran 13{,}000 pieces of model-written C++ on the host,
unsandboxed'' is the honest description, and it is a fair thing for an interviewer to
press on. The fix is well-trodden (run the binary in a container, or under
\code{bubblewrap}, or as a nobody user in a chroot) and the code is structured such that it
would be a change to \code{\_run\_one} alone.
\end{honest}

\subsection{Early stopping, and how it deflates scores}

\begin{lstlisting}[language=Python, caption={evaluator.py, lines 146-180. Note which denominator the score uses.}]
for i in range(n):
    actual = _run_one(bin_path, inputs[i], timeout)
    ok     = actual is not None and _match(actual, outputs[i])
    res["tests_run"] = i + 1
    if ok:
        passed += 1
        fails   = 0          # the counter RESETS on a pass: 3 CONSECUTIVE failures needed
    else:
        fails += 1
        if fails >= config.EARLY_STOP_FAILURES:   # = 3
            res["early_stopped"] = True
            break

res["score"] = passed / n    # <-- n is the TOTAL test count, not tests_run
\end{lstlisting}

Two things are happening, and only one is obvious:

\begin{enumerate}
  \item \textbf{Early stopping is a cost control.} Once a program has failed three tests in
  a row, it is almost certainly wrong, so the remaining tests are skipped. Over 65 problems
  $\times$ 5 individuals $\times$ 12 generations that saves a lot of wall-clock time.

  \item \textbf{The score divides by the total, not by the tests actually run.} When a run
  stops early at test 4 of 30, the score is $\mathrm{passed}/30$, not $\mathrm{passed}/4$.
  The 27 unrun tests are silently counted as failures.
\end{enumerate}

This second point is a deliberate and correct choice: assuming the unrun tests would have
failed is the conservative reading, and it keeps the score comparable between an
early-stopped evaluation and a complete one. It is worth being able to state clearly,
because it looks like a bug and is not.

\begin{honest}{Exact-match comparison is strict, and quietly wrong for some problems}
\code{\_match} normalises by stripping trailing whitespace from each line and comparing the
line lists for exact equality. That means a \emph{correct} program is scored as failing if
it prints \code{3.14159} where the expected file says \code{3.14160}, or if it picks a
different but equally valid answer to a problem with multiple correct outputs. APPS contains
both categories.

The README argues this ``deflates every prompt's score by roughly the same amount''. That
is probably true on average and is a reasonable defence, \emph{but it is not free}: it also
injects noise into the fitness landscape the GA is trying to climb, and with a population of
5 there is little averaging to damp it. A per-problem checker (float tolerance, answer-set
comparison) is the correct fix and the README says so.

A smaller note: the README refers to this logic as ``\code{\_normalise}/\code{\_compare} in
\code{evaluator.py}, line 41''. There is no function called \code{\_compare}; it is
\code{\_match}, at line 40. Minor, but it is the kind of thing an interviewer checks.
\end{honest}

\subsection{Parallelism}

\begin{lstlisting}[language=Python, caption={ga.py, lines 189-208. Five problems evaluated at once.}]
with ThreadPoolExecutor(max_workers=config.PARALLEL_EVALS) as executor:   # = 5
    futures = {executor.submit(_eval_one_problem, gen, ind, prob, i, len(problems)): i
               for i, prob in enumerate(problems)}
    for future in as_completed(futures):
        idx = futures[future]
        try:
            score, f = future.result()
            scores[idx] = score
            if f is not None:
                with lock: fails.append(f)
        except Exception as exc:
            _log(f"ind={ind.uid}  prob={idx+1}  UNEXPECTED ERROR: {exc}", level="ERROR")
            scores[idx] = -1.0
\end{lstlisting}

\begin{jargon}{Thread pool, and why threads are the right choice here}
A \textbf{thread} lets one program do several things at once. Python's threads are limited
by the Global Interpreter Lock, so they cannot speed up pure Python computation. But they
are perfectly good at \emph{waiting}, and waiting is all this workload does: waiting for a
network reply from Gemini, and waiting for \code{g++} and the test binary (both separate
processes). \code{ThreadPoolExecutor(max\_workers=5)} keeps five such waits in flight.
\code{as\_completed} yields futures as they finish, in completion order, which is why the
code keys a dictionary from future back to index rather than relying on order.
\end{jargon}

The pool size of 5 is a rate-limit compromise, not a hardware one. There is also a
deliberate jitter, \code{time.sleep(random.uniform(0, 1))} at the top of every evaluation
(\code{ga.py} line 167) and \code{uniform(0, 2)} in the holdout path, to stop five threads
from hitting the API in the same instant at the start of each batch.

The exception handler writes \code{-1.0}, the ``this is not a real score'' sentinel, rather
than 0.0. That distinction is carried carefully: \code{-1.0} means infrastructure broke,
0.0 means the model genuinely produced code that failed everything.

\begin{honest}{The sentinel is honest inside a generation and lost at the boundary}
\code{compute\_fitness} computes \code{sum(scores) / len(scores)} over the raw list, so a
\code{-1.0} sentinel is averaged in \emph{as a negative number}. An individual that hit one
infrastructure failure across 65 problems is therefore penalised by roughly $1/65 \approx
0.0154$ of fitness more than an individual whose code merely failed every test on that
problem. It is a small effect, and arguably an acceptable one (it makes flaky individuals
slightly less attractive), but it is not what a reader would assume, and the distinction the
\code{-1.0} sentinel exists to preserve is silently discarded at exactly the point it
mattered. The generation-level statistics are more careful: they filter to
\code{valid = [f for f in fits if f >= 0.0]} before computing standard deviation. So the
same sentinel is treated two different ways within thirty lines of each other.
\end{honest}

\newpage
%=============================================================================
\section{The dataset and the split}
%=============================================================================

\begin{jargon}{APPS, and Codeforces ratings}
\textbf{APPS} (Hendrycks et al., 2021) is a benchmark of programming problems scraped from
sites like Codeforces, each with a problem statement, a set of test inputs and their
expected outputs. \textbf{Codeforces rating} is a difficulty number: 800 is beginner-level,
1900 is genuinely hard. The project's subset spans 800 to 1900.
\end{jargon}

\code{dataset.py} does not scan directories. It hardcodes an explicit registry
(\code{\_SELECTED}, lines 9-58) of exactly which problem folders to load, grouped by split
and rating. Verified by running it: 300 problems load, 100 from \code{train/} and 200 from
\code{test/}.

\begin{figure}[H]
\centering
\resizebox{0.92\textwidth}{!}{
\begin{tikzpicture}[node distance=9mm]

  \node[store, minimum width=26mm] (all) {300 problems\\\code{\_SELECTED} registry};

  \node[box, right=16mm of all, minimum width=26mm, yshift=20mm] (train) {\textbf{train/} split\\100 problems\\ratings 800-1900};
  \node[box, right=16mm of all, minimum width=26mm, yshift=-16mm] (test) {\textbf{test/} split\\200 problems\\ratings 800-1400};

  \node[box, right=16mm of train, minimum width=30mm, yshift=13mm] (hard) {rating $\ge$ 1800\\\textbf{5 problems}\\always in evo pool};
  \node[box, right=16mm of train, minimum width=30mm, yshift=-13mm] (tiers) {rating $<$ 1800\\stratified quota\\per \code{\_EVOLUTION\_PER\_TIER}};

  \node[box, right=16mm of $(hard)!0.5!(tiers)$, minimum width=26mm, draw=okgreen] (evo) {\textbf{evolution pool}\\\textbf{65 problems}\\fitness is measured here};

  \node[box, right=16mm of test, minimum width=26mm, draw=warnred] (hold) {\textbf{holdout}\\\textbf{200 problems}\\touched once, at the end};

  \node[box, below=8mm of hold, minimum width=26mm] (rest) {\textbf{rest}\\35 problems\\unused};

  \draw[flow] (all) -- (train);
  \draw[flow] (all) -- (test);
  \draw[flow] (train) -- (hard);
  \draw[flow] (train) -- (tiers);
  \draw[flow] (hard) -- (evo);
  \draw[flow] (tiers) -- (evo);
  \draw[flow] (test) -- node[lbl,pos=0.5]{shuffled,\\first 200} (hold);
  \draw[flow] (test) -- (rest);

\end{tikzpicture}}
\caption{The split, verified by running \code{dataset.split}: evolution 65, holdout 200,
rest 35. The five hardest train problems are forced into the evolution pool; the remaining
60 slots are filled by a per-rating quota.}
\end{figure}

\begin{keyidea}{Why the evolution pool and the holdout must be different problems}
If you tune a prompt on the same problems you then score it on, a high score proves nothing:
the prompt may simply have absorbed the quirks of those 65 problems. This is
\textbf{overfitting}. The 200-problem holdout is never seen during evolution and is
evaluated exactly once, at the very end, on the single winning prompt. That is the number
that means something.

The split is reproducible (\code{RANDOM\_SEED = 42}) and, better, the actual chosen IDs are
written to \code{evolution\_ids.json} in the run directory, so a resumed run restores the
\emph{exact same} split from disk rather than trusting the seed to regenerate it. That is
the more robust choice: it survives a change to the splitting code.
\end{keyidea}

The stratification is deliberate. \code{\_EVOLUTION\_PER\_TIER} (lines 61-71) fixes a quota
per rating tier (7 at rating 800, 6 at 1000, and so on), so the evolution pool cannot end up
accidentally all-easy. The five rating-1800-and-above problems bypass the quota entirely
and are always included: with only five hard problems available, leaving their inclusion to
chance would risk a pool with no hard problems at all.

\begin{honest}{Three real problems with the dataset layer}
\textbf{The holdout is drawn only from \code{test/}, the evolution pool only from
\code{train/}.} Look at the \code{\_SELECTED} registry: \code{train/} spans ratings 800 to
1900, \code{test/} spans only 800 to 1400. So the evolution pool contains problems rated
1500, 1600, 1700, 1800 and 1900 that have \emph{no counterpart} in the holdout, and the
holdout's difficulty distribution is materially easier than the pool the prompts were tuned
on. The generalisation claim is therefore ``tuned on 800-1900, tested on 800-1400''. That is
not wrong, but it is not the like-for-like comparison a reader assumes, and neither the
README nor the paper mentions it.

\textbf{\code{reserved\_train\_problems} (lines 163-179) is dead code.} It is never called
from anywhere in the repository. It computes a ``reserved'' set of train problems using
\code{\_RESERVED\_PER\_TIER}, a table that exists solely to feed this unused function. It
appears to be a leftover from a planned third split.

\textbf{\code{sys.set\_int\_max\_str\_digits(1\_000\_000)}} is called at the top of
\code{load\_problems} (line 87) with no comment. It raises Python 3.11's guard against
converting enormous integers to strings, which exists to prevent a denial-of-service. It is
presumably needed because some APPS problem has a huge numeric literal in its test data, but
nothing records that, and a reader cannot tell whether it is load-bearing or superstition.
\end{honest}

\newpage
%=============================================================================
\section{The LLM layer}
%=============================================================================

\code{llm.py} is the only file that talks to Google Cloud. It handles authentication,
retries, token accounting, and the four prompt-shaped functions the GA calls
(\code{generate\_code}, \code{crossover}, \code{mutate\_inject}, \code{mutate\_delete},
\code{mutate\_rephrase}).

\subsection{Client construction: double-checked locking}

\begin{lstlisting}[language=Python, caption={llm.py, lines 87-126, condensed. One client, shared by five threads.}]
_api_client = None
_api_lock   = threading.Lock()

def _get_client():
    global _api_client
    if _api_client is not None:      # fast path: no lock taken once initialised
        return _api_client
    with _api_lock:
        if _api_client is not None:  # re-check INSIDE the lock
            return _api_client
        ...
        _api_client = genai.Client(vertexai=True, project=project,
                                   location=config.GCP_REGION, credentials=creds,
                                   http_options=_gt.HttpOptions(timeout=120_000))
    return _api_client
\end{lstlisting}

\begin{jargon}{Double-checked locking}
Five threads may call \code{\_get\_client} simultaneously at startup. Without a lock, all
five could see \code{None} and each build a client. Taking a lock on \emph{every} call would
serialise the fast path forever after. So: check without the lock (cheap, and correct once
initialised), and only if it looks uninitialised take the lock and check \emph{again}
inside it, because another thread may have won the race while you waited. This idiom is
famously subtle in some languages; in CPython it is safe here because assignment to a module
global is atomic.
\end{jargon}

\subsection{The retry state machine}

\begin{figure}[H]
\centering
\resizebox{0.95\textwidth}{!}{
\begin{tikzpicture}

  \node[box, minimum width=26mm] (try)  at ( 0.0,   0.0) {attempt $n$\\(max 10)};
  \node[dec, minimum width=18mm] (exc)  at ( 0.0,  -3.0) {exception?};

  \node[dec, minimum width=20mm] (empty) at ( 6.8, -3.0) {text empty?};
  \node[box, minimum width=28mm, draw=okgreen] (ok) at (13.2, -3.0) {\textbf{return text}\\record tokens $+$ cost};

  \node[dec, minimum width=22mm] (kind) at ( 0.0,  -6.3) {429 / 5xx /\\timeout?};
  \node[box, minimum width=34mm] (wait) at ( 0.0,  -9.6) {sleep $5 \cdot 2^{n} + U(0,3)$\\seconds, then retry};
  \node[box, minimum width=30mm, draw=warnred] (giveup) at ( 0.0, -12.9) {retries exhausted\\\textbf{return None}\\$\rightarrow$ LLMCallError};

  \draw[flow] (try) -- (exc);
  \draw[flow] (exc)   -- node[lbl, pos=0.42] {no} (empty);
  \draw[flow] (empty) -- node[lbl, pos=0.42] {no} (ok);
  \draw[flow] (exc)   -- node[lbl, pos=0.5] {yes} (kind);
  \draw[flow] (kind)  -- node[lbl, pos=0.5] {yes; and in fact \emph{no} too:\\every error is retried} (wait);
  \draw[flow] (wait)  -- node[lbl, pos=0.5] {after 10 attempts} (giveup);

  % empty response is retryable: down the right, then back into the wait
  \draw[flow] (empty.south) -- node[lbl, pos=0.28] {yes: retry} (6.8,-9.6) -- (wait.east);

  % the retry loop, routed up the left, clear of both decisions
  \draw[dflow] (wait.west) -- (-5.2,-9.6) -- node[lbl, pos=0.55] {attempt $n{+}1$} (-5.2,0.0) -- (try.west);

\end{tikzpicture}}
\caption{\code{llm.\_call\_api}, lines 142-220. Exponential backoff with jitter, up to 10
attempts, covering both exceptions and the empty-response case.}
\end{figure}

\begin{jargon}{Exponential backoff with jitter}
When a server says ``too many requests'', retrying immediately makes it worse. So wait, and
double the wait each time: 5s, 10s, 20s, 40s and so on ($5 \cdot 2^{n}$). The random extra
(\code{random.uniform(0, 3)}, the \emph{jitter}) stops five threads that were rate-limited
at the same instant from retrying at the same instant and colliding again. Without jitter,
parallel clients synchronise into waves.
\end{jargon}

The empty-response retry deserves note: a Gemini call can succeed at the HTTP level and
return no text (the model spent its whole token budget on internal reasoning, or a safety
filter fired). The code treats that as retryable, which is right. Critically, it records the
cost \emph{before} checking whether the text is empty (lines 178-188), so an empty response
still gets billed and still gets counted. That is the honest accounting choice: Google
charges for it regardless.

\begin{honest}{The backoff can wait 85 minutes, and every error is retried}
At attempt 9 the wait is $5 \cdot 2^{9} = 2560$ seconds, over 42 minutes; the cumulative
worst case across 10 attempts is roughly 85 minutes for a \emph{single} call, with no cap.
Worse, the final \code{else} branch (lines 214-217) retries even errors that were explicitly
classified as neither rate-limit, server error, nor timeout: an authentication failure or a
malformed request will be retried ten times with full backoff before failing. A capped
backoff (\code{min(wait, 60)}) and a non-retryable error class are both one-line changes.
\end{honest}

\begin{honest}{The configured temperatures are dead constants}
\code{config.py} defines \code{TARGET\_TEMPERATURE = 0.0} and
\code{OPTIMIZER\_TEMPERATURE = 0.7}. Verified by grep: \textbf{neither is read anywhere in
the codebase}. \code{\_call\_api} hardcodes the same values inline:
\begin{lstlisting}[language=Python]
is_fit = (phase == "fitness")
config=_gt.GenerateContentConfig(
    temperature=0.0 if is_fit else 0.7,                       # not config.TARGET_TEMPERATURE
    thinking_config=_gt.ThinkingConfig(thinking_budget=512 if is_fit else 2048),
)
\end{lstlisting}
The values agree today, so nothing is broken. But editing \code{config.py} to change a
temperature would have \emph{no effect}, silently, which is the worst kind of configuration
bug. The thinking budgets (512 and 2048) are hardcoded with no config constant at all.
\end{honest}

\begin{honest}{Two pricing tables, one imported and never used}
\code{llm.py} line 10 reads \code{from cost\_tracker import tracker, PRICING}, and then
never uses \code{PRICING}. Instead it defines its \emph{own} \code{\_pricing} dict at line
129 with overlapping but not identical contents, and \code{\_cost} uses that one. So the
project carries two pricing tables that can silently drift apart. Since \code{llm.py} always
passes an explicit \code{cost\_usd} into \code{tracker.record}, \code{cost\_tracker}'s own
\code{PRICING} is only reachable through the \code{CallRecord} properties when
\code{cost\_usd} is \code{None}, which never happens in this codebase. The whole
\code{PRICING} table plus the three \code{CallRecord} cost properties are effectively dead.
\end{honest}

\begin{honest}{run\_fitness\_batch is a stub}
\code{llm.py} lines 242-243 define \code{run\_fitness\_batch(requests)} whose entire body
raises \code{NotImplementedError("Batch API not used: realtime only.")}. Nothing calls it.
It is a fossil of an abandoned plan to use Vertex AI's batch endpoint, which would have been
roughly half the price. Similarly \code{evaluator.run\_tests} (line 191) is a thin wrapper
around \code{run\_tests\_verbose} that nothing calls.
\end{honest}

\newpage
%=============================================================================
\section{Cost tracking}
%=============================================================================

\begin{keyidea}{Why a course project needed a cost subsystem}
Every fitness evaluation is a paid API call. One run is roughly 4{,}000 calls over 4 to 6
hours, unattended. Without live accounting you start a run, go to sleep, and find out what
it cost from the bill. \code{cost\_tracker.py} writes an updated JSON report after
\emph{every single call}, so a run can be watched from another terminal and killed if it
drifts. The three committed runs cost \$5.99, \$6.16 and \$4.90, all verified from the
committed reports in \code{Costs/}.
\end{keyidea}

The tracker is a single global instance (\code{tracker = CostTracker()}, line 191) with a
lock around all mutation, accumulating totals by phase (\code{fitness}, \code{crossover},
\code{mutate\_inject}, ...) and by model. It writes three artifacts: an append-only
\code{calls.jsonl} with one line per call, a live summary at \code{costs\_live.json} in the
run directory, and a second live copy under \code{Costs/}.

\code{restore\_from\_jsonl} (lines 77-112) is what makes \code{--resume} honest about money:
it replays the entire \code{calls.jsonl} to rebuild the running total, so a resumed run
reports cumulative cost across both halves rather than restarting the counter at zero.

One subtlety handled correctly, in \code{llm.py} lines 172-175:

\begin{lstlisting}[language=Python]
in_tok  = getattr(usage, "prompt_token_count",     0) or 0
out_tok = getattr(usage, "candidates_token_count", 0) or 0
think   = getattr(usage, "thoughts_token_count",   0) or 0
out_tok += think     # thinking tokens are billed at the OUTPUT rate
\end{lstlisting}

\begin{jargon}{Thinking tokens}
Gemini 2.5 models can do hidden internal reasoning before answering. Those tokens are not
in the visible response but they are billed, at the output rate. Folding \code{think} into
\code{out\_tok} before costing is the difference between a cost estimate that matches the
bill and one that badly understates it, especially for the operator calls which are given a
2048-token thinking budget. The \code{or 0} guards against the API returning \code{None}
rather than omitting the field, which \code{getattr}'s default alone would not catch.
\end{jargon}

\begin{honest}{Three blemishes on an otherwise solid subsystem}
\textbf{Stale pricing rows.} \code{PRICING} still carries
\code{claude-haiku-4-5@20251001}, \code{claude-haiku-4-5-20251001} and
\code{claude-sonnet-4-6} rows from an earlier iteration that used a different backend. The
project has used only Gemini for the whole of its committed history. The README already
admits this one.

\textbf{A write to disk on every call, twice.} \code{record()} calls \code{\_write\_summary}
for \emph{both} live paths while holding the lock, so every one of roughly 4{,}000 calls per
run rewrites two JSON files synchronously with the global lock held. It is harmless at this
scale (the calls take seconds, the write takes microseconds) but it is a real serialisation
point, and it means a crash mid-write can leave a truncated JSON file.

\textbf{Rounding accumulates inconsistently.} Per-bucket totals are re-rounded to 6 decimal
places on every update (\code{round(bucket[key]["total\_cost\_usd"] + call\_cost, 6)}) while
the grand total is not rounded until it is written. The two can therefore disagree in the
last decimal. Immaterial in dollars; the kind of thing worth noticing.
\end{honest}

\newpage
%=============================================================================
\section{Orchestration, checkpoints and resume}
%=============================================================================

\code{main.py} is the entry point. It runs pre-flight checks, builds the split, wires the
loggers, runs the GA, then evaluates the winner on the holdout.

\subsection{Pre-flight}

The pre-flight sequence (lines 200-225) checks, in order: the service-account file exists;
\code{google-genai} and \code{google-auth} import; the service account parses and its
\code{project\_id} and \code{client\_email} are printed; both APPS directories exist; the
problems load; and enough problems survive to fill both splits. Each failure exits with a
specific message. For a program that is about to spend six hours and six dollars, failing in
the first two seconds when the credential is missing is exactly right.

\subsection{Signal handling}

\begin{lstlisting}[language=Python, caption={main.py, lines 337-343. Ctrl-C does not lose the cost record.}]
def _on_signal(signum, frame):
    print(f"\n[MAIN] signal {signum}: flushing costs and exiting ...", flush=True)
    tracker.flush()
    sys.exit(1)

signal.signal(signal.SIGTERM, _on_signal)
signal.signal(signal.SIGINT,  _on_signal)
\end{lstlisting}

\begin{honest}{The signal handler cannot do what it looks like it does}
It flushes the cost report, which is genuinely useful. But it does \emph{not} save the
current population, so pressing Ctrl-C mid-generation discards up to a full generation of
work: potentially an hour of wall clock and a dollar of API spend. The resume path can only
restart from the last completed generation's checkpoint.

There is also a subtler issue: \code{sys.exit()} in a signal handler raises
\code{SystemExit} in the \emph{main} thread only. With a \code{ThreadPoolExecutor} whose
worker threads are mid-call, the interpreter will wait on the pool's non-daemon threads at
shutdown, so Ctrl-C can appear to hang for up to the 120 second API timeout. Not a
correctness bug, but it explains behaviour that would otherwise look like a freeze.
\end{honest}

\subsection{Checkpointing and resume}

\begin{figure}[H]
\centering
\resizebox{0.96\textwidth}{!}{
\begin{tikzpicture}[node distance=8mm]

  \node[box, minimum width=26mm] (start) {\code{python main.py --resume}};
  \node[box, right=12mm of start, minimum width=26mm] (find) {\code{\_find\_latest\_run}\\newest \code{run\_*} by mtime};
  \node[dec, right=12mm of find, minimum width=20mm] (has) {generation\_*.json\\present?};

  \node[box, right=12mm of has, minimum width=28mm, yshift=15mm] (loadgen) {\code{\_load\_resume\_state}\\population + history\\+ RNG state};
  \node[box, right=12mm of has, minimum width=28mm, yshift=-15mm] (loadseed) {seed checkpoints only\\$\rightarrow$ resume gen 0\\(re-use scored seeds)};

  \node[box, right=13mm of loadgen, minimum width=26mm] (ids) {restore split from\\\code{evolution\_ids.json}};
  \node[box, below=10mm of ids, minimum width=26mm] (cost) {\code{restore\_from\_jsonl}\\replay \code{calls.jsonl}\\$\rightarrow$ cumulative cost};
  \node[box, below=10mm of cost, minimum width=26mm, draw=okgreen] (go) {continue from\\gen $+$ 1};

  \draw[flow] (start) -- (find);
  \draw[flow] (find) -- (has);
  \draw[flow] (has) -- node[lbl,pos=0.4]{yes} (loadgen);
  \draw[flow] (has) -- node[lbl,pos=0.4]{no} (loadseed);
  \draw[flow] (loadgen) -- (ids);
  \draw[flow] (loadseed.east) to[bend right=24] (ids.south west);
  \draw[flow] (ids) -- (cost);
  \draw[flow] (cost) -- (go);

\end{tikzpicture}}
\caption{The resume path. Three separate pieces of state are restored: the population, the
random number generator, and the cumulative cost. Restoring the RNG is what makes a resumed
run take the same decisions the uninterrupted run would have.}
\end{figure}

\begin{keyidea}{Why the RNG state is checkpointed, and why that is unusually careful}
Seeding a generator with \code{Random(42)} makes a run reproducible \emph{from the start}.
It does nothing for a run resumed at generation 7: a fresh \code{Random(42)} would replay
the coin flips from generation 1, so the resumed run would take different decisions than the
uninterrupted one would have.

So \code{ga.rng\_state\_to\_dict} serialises Python's full Mersenne Twister internal state
(a 625-element tuple) into JSON, and \code{rng\_state\_from\_dict} restores it. The resumed
run's next coin flip is precisely the one the original run would have made. Most student
projects do not do this, and most would not notice it was missing.
\end{keyidea}

Seed evaluation gets its own finer-grained checkpointing: each of the five seed prompts is
scored and immediately written to \code{seed\_NN\_uid.json} (\code{ga.py}, lines 378-387).
Generation 0 is the most expensive single generation, since all five individuals are new; a
crash during it would otherwise waste the lot. On resume, already-scored seeds are restored
and skipped.

\begin{honest}{The final generation is never checkpointed, and it breaks two things}
\code{save\_checkpoint} is only ever called from \code{\_on\_gen}, which \code{ga.run}
invokes at the \emph{top} of the loop, recording generation \code{gen-1}. When the loop
finishes its last iteration, it exits and returns the winner without ever writing the final
population.

This is not a theoretical concern. Verified against all three committed runs:

\begin{center}
\small
\begin{tabular}{lccc}
\toprule
run & checkpoints & gens in \code{history.json} & winner in a checkpoint? \\
\midrule
\code{run\_20260429\_210411} (tourn.) & 12 (00 to 11) & 13 (0 to 12) & \textbf{no} \\
\code{run\_20260430\_014526} (roulette) & 12 (00 to 11) & 13 (0 to 12) & \textbf{no} \\
\code{run\_20260430\_014530} (elitism) & 12 (00 to 11) & 13 (0 to 12) & yes \\
\bottomrule
\end{tabular}
\end{center}

Two concrete consequences:

\textbf{1. Resume silently rewinds.} A \code{--resume} after a completed 12-generation run
restarts from generation 11, discarding generation 12 entirely.

\textbf{2. It is the root cause of the broken genealogy} (Section~\ref{sec:genealogy}). The
winning individual of the tournament and roulette runs exists only in \code{summary.json}
and \code{best\_prompt.txt}; it is in no checkpoint, so its ancestry cannot be traced. The
elitism run escapes only by luck: its winner (\code{822eea31}) was an elite from generation
6, carried forward, and so it appears in earlier checkpoints.

The fix is one line: call \code{on\_generation} once more after the loop.
\end{honest}

\subsection{The holdout evaluation}

After the GA returns, \code{evaluate\_on\_holdout} (lines 154-175) scores the single winning
prompt against all 200 holdout problems, in the same 5-way thread pool. Note the
difference from \code{compute\_fitness}: here an infrastructure failure is clamped to 0.0
(\code{scores[idx] = entry["score"] if entry["score"] >= 0.0 else 0.0}), rather than
averaged in as \code{-1.0}. That is the third distinct treatment of the sentinel in the
codebase, and this one is the most defensible.

\newpage
%=============================================================================
\section{The analysis scripts}
%=============================================================================

\subsection{run\_baseline\_eval.py}

\begingroup\sloppy
Scores the tournament run's evolved prompt against the strongest hand-written seed
(\code{SEED\_PROMPTS[0]}) on all 200 test-split problems, writing the results into
\code{comparison.json}. It has a nice resumability property: it reads any existing JSONL
first, keys by \code{problem\_id}, and only evaluates what is missing, so an interrupted
comparison can be restarted without re-paying for completed problems.
\par\endgroup

\begin{honest}{The script cannot reproduce three of the four result sets it is credited with}
The constants at the top are hardcoded:
\begin{lstlisting}[language=Python]
BEST_RUN_DIR   = config.RESULTS_DIR / "run_20260429_210411"      # the tournament run only
EVOLVED_JSONL  = OUT_DIR / "evolved_tournament_test_evaluations.jsonl"
EVOLVED_PROMPT = (BEST_RUN_DIR / "best_prompt.txt").read_text().strip()   # at IMPORT time
\end{lstlisting}
So the committed script produces exactly two files: \code{evolved\_tournament\_...} and
\code{baseline\_...}. But \code{results/comparison/} also contains
\code{evolved\_roulette\_test\_evaluations.jsonl} and
\code{evolved\_elitism\_test\_evaluations.jsonl}, and
\code{visualize\_test\_results.py} \emph{requires} all four and exits if any is missing.

The roulette and elitism evaluation logs were therefore produced by hand-editing these
constants between runs, and \textbf{that variant of the script was never committed}. The
headline four-way comparison in the README and the paper cannot be regenerated from this
repository without editing source. There is no scheme parameter and no CLI flag; adding one
is a fifteen-minute change that would close the gap.

Second, smaller: \code{EVOLVED\_PROMPT} is read at \emph{module import time}, not inside
\code{main()}. Importing this module anywhere without that exact run directory present
raises immediately, before any error handling can produce a useful message.
\end{honest}

\begin{honest}{comparison.json and the figures compute the mean two different ways}
\code{build\_comparison} computes \code{mean\_score} over \emph{valid} entries only
(\code{[r["score"] for r in results if r.get("score", -1) >= 0]}, divided by
\code{len(valid)}). \code{visualize\_test\_results.py} divides by a hardcoded 200
(\code{sum(v["score"] for v in t.values()) / 200}), which would count a \code{-1.0} sentinel
as a \emph{negative}.

The two agree today, and I verified why: all four evaluation logs contain 200 entries with
zero infrastructure failures, so \code{len(valid) == 200} and the sentinel never appears.
The agreement is luck, not design. A single \code{infra\_failure} in any log would make the
JSON and the figures disagree, and the figures would be wrong.
\end{honest}

\subsection{visualize\_test\_results.py}

Regenerates the five comparison figures from the committed JSONL logs. It needs no API
access, and I confirmed it works: running it reproduced all five PNGs and printed exactly
the numbers the README claims (Section~\ref{sec:verified}). It opens with
\code{matplotlib.use("Agg")} before importing \code{pyplot}, selecting the non-interactive
backend so it runs on a headless machine, and asserts up front that all four logs have 200
entries with matching problem IDs. Those assertions are the right instinct: a silent
mismatch between prompt runs would produce a plausible-looking and completely wrong chart.

\newpage
%=============================================================================
\section{What the committed results actually say}
\label{sec:verified}
%=============================================================================

Everything in this section was re-derived from the committed artifacts, not taken from the
README.

\subsection{The headline: the GA lost}

\begin{table}[H]
\centering
\begin{tabular}{lcc}
\toprule
Prompt & Mean pass rate (200 test problems) & vs baseline \\
\midrule
\textbf{Baseline} (best hand-written seed) & \textbf{0.7900} & \\
Tournament evolved & 0.7857 & $-0.0043$ \\
Roulette evolved   & 0.7717 & $-0.0183$ \\
Elitism evolved    & 0.7713 & $-0.0187$ \\
\bottomrule
\end{tabular}
\caption{Verified by recomputing from the four \code{*\_test\_evaluations.jsonl} files, and
independently reproduced by re-running \code{visualize\_test\_results.py}. All four logs
contain exactly 200 entries and zero infrastructure failures. \textbf{No evolved prompt beat
the hand-written baseline.}}
\end{table}

Per-problem, tournament vs baseline: \textbf{26 wins, 22 losses, 152 ties} (verified by
recomputing the deltas in \code{comparison.json}). Tournament was the only scheme with a
positive head-to-head record, and 152 ties out of 200 tells its own story: on three quarters
of the problems the prompt made no difference whatsoever.

The outcome distributions are strikingly flat. Baseline: 142 perfect, 44 partial, 12 compile
failures, 2 total failures. Tournament evolved: 142 perfect, 44 partial, 10 compile
failures, 4 total failures. Evolution moved two problems from ``compile failure'' to ``fails
all tests''. That is the entire measured effect of the whole exercise on this metric.

\begin{keyidea}{The honest headline, and why it is still worth defending}
The project's result is \textbf{negative}: a genetic algorithm, given \$17 and 13{,}000
model calls, did not find an instruction prompt better than the one a human wrote by hand in
an afternoon.

That is a real, publishable finding, and the machinery that produced it is sound. The value
of the project is that it is \emph{capable of proving its own hypothesis wrong}, which most
prompt-optimisation work is not, because most of it has no objective fitness function and no
held-out test set. Defend the method, report the result.
\end{keyidea}

\subsection{Cost, verified}

\begin{table}[H]
\centering
\begin{tabular}{lrr}
\toprule
Run & LLM calls & Cost (USD) \\
\midrule
\code{run\_20260429\_210411} (tournament) & 4{,}734 & \$5.99 \\
\code{run\_20260430\_014526} (roulette)   & 4{,}529 & \$6.16 \\
\code{run\_20260430\_014530} (elitism)    & 3{,}739 & \$4.90 \\
\midrule
\textbf{Total} & \textbf{13{,}002} & \textbf{\$17.05} \\
\bottomrule
\end{tabular}
\caption{Read directly from \code{Costs/cost\_run\_*.json}. This matches the README's
``roughly 13,000 LLM calls and about \$17'' exactly. Note the roulette run cost \emph{more}
than the tournament run despite 205 fewer calls: roulette's operator calls to the expensive
Gemini 2.5 Pro must have run longer or produced more thinking tokens.}
\end{table}

\subsection{The finding the README misses: evolution went backwards}
\label{sec:regression}

This is the most important thing in this document, and neither the README nor the paper says
it. It is derived entirely from the committed \code{history.json} of each run.

\begin{figure}[H]
\centering
\begin{tikzpicture}
\begin{axis}[
  width=0.92\textwidth, height=7cm,
  xlabel={generation}, ylabel={best fitness in population},
  xmin=-0.4, xmax=12.4, ymin=0.62, ymax=0.76,
  grid=both, grid style={linegrey!50},
  legend style={font=\scriptsize, fill=white, draw=linegrey,
                at={(0.5,-0.20)}, anchor=north, legend columns=3, column sep=1.2em},
  label style={font=\small}, tick label style={font=\scriptsize},
  every axis plot/.append style={thick, mark size=1.6pt},
]
\addplot[color=inkblue, mark=*] coordinates {
 (0,0.6909) (1,0.6559) (2,0.6641) (3,0.6717) (4,0.6812) (5,0.7181)
 (6,0.6654) (7,0.7063) (8,0.6918) (9,0.7387) (10,0.6854) (11,0.6998) (12,0.6531)};
\addlegendentry{tournament (elitism 0)}
\addplot[color=okgreen, mark=square*] coordinates {
 (0,0.7380) (1,0.6961) (2,0.7116) (3,0.7032) (4,0.6856) (5,0.7015)
 (6,0.6648) (7,0.6744) (8,0.7084) (9,0.6484) (10,0.7034) (11,0.7005) (12,0.6619)};
\addlegendentry{roulette (elitism 0)}
\addplot[color=warnred, mark=triangle*] coordinates {
 (0,0.7338) (1,0.7338) (2,0.7338) (3,0.7338) (4,0.7338) (5,0.7338)
 (6,0.7346) (7,0.7346) (8,0.7346) (9,0.7346) (10,0.7346) (11,0.7346) (12,0.7346)};
\addlegendentry{elitism (elitism 1)}
\end{axis}
\end{tikzpicture}
\caption{Best fitness per generation, plotted from the committed \code{history.json} of each
run. The tournament and roulette curves are \emph{not} monotonic: they wander, and both
\textbf{end below where they started}. Only the elitism run is non-decreasing, which is
exactly what an elite guarantees.}
\end{figure}

\begin{honest}{The README's ``evolution reliably improved fitness'' is not what the data shows}
The README states: ``evolution reliably improved fitness on the 65-problem training pool
(best-of-run fitness 0.65 to 0.73 depending on scheme)''. Those three numbers are the
\emph{final generation's} best. Compare them with generation 0, the hand-written seeds:

\begin{center}
\begin{tabular}{lccc}
\toprule
run & gen 0 best (seeds) & gen 12 best (final) & change \\
\midrule
tournament & 0.6909 & 0.6531 & \textbf{$-0.0378$ (worse)} \\
roulette   & 0.7380 & 0.6619 & \textbf{$-0.0761$ (worse)} \\
elitism    & 0.7338 & 0.7346 & $+0.0008$ (flat) \\
\bottomrule
\end{tabular}
\end{center}

\textbf{In two of three runs the GA ended with a worse prompt than the seeds it started
from.} The third improved by 0.0008, which is noise. The roulette run's peak fitness was its
generation 0. The tournament run peaked at 0.7387 at generation 9 and then \emph{lost that
individual}, finishing 0.086 below its own peak.

The mechanism is not mysterious and is worth being able to explain: with \code{elitism = 0},
the entire population is replaced every generation, so nothing protects the best individual.
A good prompt found at generation 9 is gone by generation 10 unless it happens to be
selected \emph{and} survives crossover and mutation intact. With a population of 5 and an
LLM rewriting the text at every step, it usually does not.

The paper's claim (line 516) that ``on the 65-problem training subset, the GA consistently
produced prompts that outperformed the best'' seed is contradicted by the repository's own
\code{history.json}. So is the Discussion's claim (line 505) that ``tournament selection
with elitism delivered the most reliable training improvement due to its strict,
non-decreasing best-fitness guarantee'': that guarantee is real, but the run labelled
\emph{tournament} had elitism 0 and is visibly not non-decreasing.

\textbf{This is the single most important thing to be ready to discuss.} The defensible
position is a good one: ``the run that protected its best individual held its ground; the
two that did not, regressed. That is textbook GA behaviour and my data shows it cleanly. My
write-up reported the final generation's best as though it were the run's achievement, and
that was wrong.''
\end{honest}

\begin{honest}{The recorded elitism\_count is wrong in all three runs}
\code{main.py} writes the run's \code{config.json} using the \emph{module constant}:
\begin{lstlisting}[language=Python]
(run_dir / "config.json").write_text(json.dumps({
    ...
    "elitism_count": config.ELITISM_COUNT,     # NOT the value actually passed to ga.run
\end{lstlisting}
but the value actually used comes from \code{main}'s own argument:
\code{elitism\_count = 1 if args.scheme == "elitism" else None}. So:

\begin{center}
\begin{tabular}{llcc}
\toprule
run & scheme & \code{config.json} records & actually used \\
\midrule
\code{run\_20260429\_210411} & (none: predates \code{--scheme}) & 2 & 0 (inferred) \\
\code{run\_20260430\_014526} & roulette & 0 & 0 \\
\code{run\_20260430\_014530} & elitism  & 0 & \textbf{1} \\
\bottomrule
\end{tabular}
\end{center}

The elitism run's own config file says it used no elitism. The tournament run's says it used
\emph{two} elites.

The ``0 (inferred)'' entry needs stating precisely. What is \emph{verified}: that run's
\code{config.json} records \code{elitism\_count: 2}, and its \code{summary.json} records
\code{selection\_scheme: null}, meaning it predates the \code{--scheme} flag. What is
\emph{inferred}: elitism 2 cannot have been in force, because an elite mathematically
guarantees best fitness never decreases, and that run's best fitness decreases repeatedly
(Figure above). So either \code{config.ELITISM\_COUNT} was 0 at run time and the 2 was
written by a later edit, or the code changed mid-run. All three runs record
\code{resumed: true}, so a code change between the original run and its resume is a live
possibility. \textbf{The artifacts do not record enough to settle it.}

The lesson generalises: a run directory should record the values a run \emph{actually used},
resolved at call time, not a snapshot of module constants taken before the arguments were
applied. As it stands the committed \code{config.json} files cannot be trusted, and the
provenance of the flagship tournament run cannot be fully reconstructed.
\end{honest}

\subsection{The truncated winner}
\label{sec:truncation}

\begin{lstlisting}[caption={results/run\_20260429\_210411/best\_prompt.txt, in its entirety (160 characters).}]
First, approach the problem systematically: analyze the input-output relationship to
identify the core algorithm, paying close attention to variable types, edge
\end{lstlisting}

It stops mid-sentence, at ``edge''. Presumably ``edge cases''. This is the prompt that won
the tournament run and that \code{run\_baseline\_eval.py} scored at 0.7857 against the
baseline's 0.7900.

The cause is \code{API\_FALLBACK\_MAX\_TOKENS\_OPS = 700} in \code{config.py}, capping every
operator's output. The model was cut off mid-word, \code{res.strip()} was truthy, and the
fragment entered the population. All three winners show it:

\begin{center}
\begin{tabular}{lrl}
\toprule
run & length & ends with \\
\midrule
tournament & 160 chars & ``\code{...variable types, edge}'' \\
roulette   & 120 chars & ``\code{...first perform a}'' \\
elitism    & 132 chars & ``\code{...include main(),}'' \\
\bottomrule
\end{tabular}
\end{center}

\begin{honest}{Every single winning prompt is a truncated fragment}
This is worse than the README's framing of it as a one-run artifact. \textbf{All three}
runs' winners are cut off mid-sentence, and all three are between 120 and 160 characters:
far shorter than the roughly 500-character seeds they descend from.

That reframes the whole result. The comparison is not ``evolved prompt vs hand-written
prompt''. It is ``a truncated 160-character fragment vs a complete 500-character
hand-written prompt'', and the fragment lost by 0.0043. One could argue the GA was doing
something interesting (short prompts scoring within half a point of a careful long one is
not nothing), but what it \emph{was not} doing is the experiment as described.

The 700-token cap should never have bound at all: these prompts are 30 to 40 tokens. The
likely explanation, and this is \textbf{inferred}, not verified, is that
\code{max\_output\_tokens} is being consumed by Gemini 2.5 Pro's 2048-token thinking budget
against a 700-token output ceiling, so the visible output is truncated after the model has
already spent its allowance reasoning. That is consistent with the thinking budget being set
to 2048 for operator calls while the output cap is 700, but I could not verify it without
API access.

Either way the defect is the missing completeness check flagged in
Section~\ref{sec:operators}: an operator's output should be rejected if it does not end at
a sentence boundary, or if the API reports a length-based finish reason. The API's
\code{finish\_reason} field is available on the response and the code never looks at it.
\end{honest}

\subsection{The broken genealogy}
\label{sec:genealogy}

The README highlights ``full genealogy logging (every individual records its parents,
operators, and intermediate prompt states)'' as a headline engineering feature, and the
paper's abstract goes further: ``Genealogy analysis reveals which instructional strategies,
edge-case emphasis, pseudocode planning, output formatting, survive selection pressure and
generalize across problem distributions.''

The underlying data capture is real: \code{parent\_a}, \code{parent\_b}, \code{operators}
and \code{intermediate\_steps} are populated and are written into every
\code{generation\_NN.json}. But the analysis built on it did not work:

\begin{lstlisting}[caption={results/run\_20260429\_210411/genealogy\_tree.json, in its entirety.}]
{
  "uid": "efd4a762",
  "note": "not found in checkpoints"
}
\end{lstlisting}

\begin{lstlisting}[caption={The corresponding genealogy\_report.txt, excerpted.}]
Best uid  : efd4a762
Lineage depth (parent_a chain): 0 steps
Total individuals in registry: 60
\end{lstlisting}

\begin{honest}{The genealogy analysis found nothing, and the script that produced it is not in the repo}
Two separate problems.

\textbf{It failed.} The tournament and roulette reports both show ``Lineage depth: 0
steps'' and a tree of \code{"not found in checkpoints"}. Only the elitism run produced a
real lineage (depth 2). The root cause is established in the checkpointing box above: the
final generation is never written to disk, so the winners of the two elitism-free runs exist
in no checkpoint and their ancestry cannot be traced. The feature is not conceptually broken;
it is defeated by a one-line omission elsewhere.

\textbf{The script does not exist.} Verified by grep: \textbf{no Python file in the
repository contains the string ``genealogy''}. \code{genealogy\_report.txt} and
\code{genealogy\_tree.json} are committed \emph{outputs} of a script that was never
committed. So the repository's headline analysis feature cannot be re-run, inspected, or
fixed by anyone reading it, including its author in six months.

Consequently the paper's abstract claim that genealogy analysis ``reveals which
instructional strategies survive selection pressure'' is not supported by anything in this
repository. The committed genealogy analysis reveals a depth of 0 for the flagship run. The
raw data to do the analysis properly is present in the checkpoints; the analysis is not.
\end{honest}

\newpage
%=============================================================================
\section{Where the documentation and the code disagree}
%=============================================================================

Collected in one place, because these are the questions an interviewer asks.

\subsection{The license contradiction}

\begin{honest}{The README says MIT. The LICENSE file is not MIT.}
README line 79: ``Code is MIT licensed (see LICENSE).''

\code{LICENSE} line 3: ``\code{\# PolyForm Strict License 1.0.0}''.

These are close to opposites. MIT permits reuse, modification and redistribution;
PolyForm Strict 1.0.0 permits \emph{viewing only} and forbids any use, reuse or
redistribution. The git history explains how it happened: commit \code{38c7f8b}, ``Switch
license to PolyForm Strict 1.0.0 (public to view, no reuse or redistribution)'', changed
the file and never updated the README sentence pointing at it.

This is the most consequential documentation bug in the repository, because it is the one
with legal meaning: a reader who trusts the README could reasonably believe they had been
granted MIT rights that the actual license withholds. It is a one-line fix and it should be
made before anyone else reads the repo.
\end{honest}

\subsection{The g++ failure mode is described backwards}

\begin{honest}{The README says it fails loudly. It scores zero silently.}
README line 83 (Challenges): ``no \code{g++} on PATH means no fitness, and the fitness path
fails with a plain `g++ not found' rather than silently scoring zero.''

Verified on this machine, which has \code{gcc-13} but no \code{g++}:
\begin{lstlisting}
>>> evaluator.run_tests_verbose('int main(){return 0;}', ['1'], ['1'], 2)
reason: compile_fail | compile error: 'g++ not found on PATH' | score: 0.0
\end{lstlisting}

It returns \code{score = 0.0} and \code{reason = "compile\_fail"}, and the GA carries on
perfectly happily. The string ``g++ not found on PATH'' appears only inside the per-problem
JSON log, indistinguishable at the fitness level from a genuine compile error. A run on a
machine without \code{g++} would complete all 12 generations, spend the full \$6 on API
calls, assign every individual a fitness of exactly 0.0, and report a winner chosen at
random. \textbf{The README describes the opposite of what the code does.}

Note that \code{RUNNING.md} gets this right: its failure table says ``\code{g++ not found on
PATH} in results $\rightarrow$ Install g++; every candidate scores 0 without it''. So the
repository contradicts itself, and the accurate description is in the less prominent file.

The fix belongs in \code{main.py}'s pre-flight, which already checks the credential and the
dataset but not the compiler: a \code{shutil.which("g++")} check, or a trial compile of a
hello-world, would turn a six-hour \$6 no-op into a two-second error.
\end{honest}

\subsection{The paper's abstract claims comparisons that were never run}

\begin{honest}{Chain-of-thought and random search do not exist in this codebase}
The paper's abstract states: ``Preliminary results demonstrate that evolved prompts
outperform hand-crafted baselines, chain-of-thought prompts, and random search, achieving
meaningful improvements in pass@1 scores.''

Three claims, and the committed repository supports none of them:
\begin{itemize}
  \item \textbf{``outperform hand-crafted baselines''}: contradicted. Evolved 0.7857 vs
  baseline 0.7900. The README already flags this one honestly.
  \item \textbf{``chain-of-thought prompts''}: verified by grep, \textbf{there is no
  chain-of-thought comparison anywhere in the codebase}. No script, no result file, no
  mention.
  \item \textbf{``random search''}: likewise, \textbf{no random-search baseline exists}. It
  is not implemented and no artifact records one.
\end{itemize}

The only comparison the repository can support is evolved vs one hand-written seed, and it
went the wrong way. The word ``preliminary'' does some work, but not this much. If the paper
is ever shown to anyone, this abstract is the first thing they read and the easiest thing to
falsify.
\end{honest}

\subsection{The paper describes a configuration that does not exist}

The paper (Experimental Configurations) describes a ``full-scale setup: population size 15,
training dataset 100 problems, 20 generations'' alongside the ``optimised setup'' of
population 5 / 65 problems / 12 generations that the code actually implements.

There is no committed artifact from a population-15, 20-generation run, and \code{config.py}
has no such setting. One tantalising trace: \code{seeds.py} contains exactly \textbf{15}
seed prompts, of which \code{main.py} uses only the first five
(\code{SEED\_PROMPTS[:config.POPULATION\_SIZE]}). The remaining ten are almost certainly the
population-15 setup's seeds, left in place when the configuration shrank. This is
\textbf{inferred} from the coincidence of the number 15 and the slicing, not confirmed
anywhere in the source.

Note this also makes the README's ``starts from 5 hand-written seeds in \code{seeds.py}''
imprecise: there are 15 in the file, and 5 are used.

\subsection{corrections.tex: a patch file, partially applied}

\code{corrections.tex} is a loose file at the repo root containing blocks of LaTeX marked
``REPLACE: ... (was lines 263-277)'' to be pasted into the paper. Checking whether they
landed: the ``Elitism (top-1)'' correction \emph{is} in the paper (line 272), so the file
was partially applied. The survival-scheme wording differs from the corrected version, and
the Table 2 corrections (adding the Optimizer LLM row and the temperatures) are not present
in the paper's parameter table.

So \code{corrections.tex} is a half-applied patch left in the repository root with no
indication of which parts landed. It should either be applied and deleted, or moved
somewhere that marks it as history. As a public-facing artifact it reads as an unfinished
edit.

\subsection{The run directory naming drifted}

Both the code and \code{RUNNING.md} say a run directory carries its scheme in the name. The
committed ones do not:

\begin{lstlisting}
main.py line 292 creates:   run_{timestamp}_{args.scheme}
RUNNING.md documents:       results/run_<timestamp>_<scheme>/

but the three committed run directories are bare timestamps:
    results/run_20260429_210411      (no scheme suffix)
    results/run_20260430_014526      (no scheme suffix)
    results/run_20260430_014530      (no scheme suffix)
\end{lstlisting}

The committed runs therefore predate the current naming code, which is consistent with the
tournament run's \code{summary.json} recording \code{selection\_scheme: null}. Harmless in
itself, but it is direct evidence that \textbf{the committed results were not produced by
the committed code}, which is the thing to be careful about: every number in this repository
comes from a version of the program that no longer exists in it.

\subsection{Minor}

\code{main.py} line 222 prints ``loading 305 pre-selected Codeforces problems'' and then
loads 300 (verified). The README already notes this one.

\newpage
%=============================================================================
\section{Verification record}
%=============================================================================

What was actually executed while writing this document, on a machine with Python 3.11+,
matplotlib and numpy, but \textbf{no \code{g++}} and \textbf{no GCP credentials}:

\begin{table}[H]
\centering
\small
\begin{tabular}{p{6.6cm}p{8.2cm}}
\toprule
\textbf{Check} & \textbf{Result} \\
\midrule
All 10 modules byte-compile & Pass \\
\code{dataset.load\_problems()} & Loads \textbf{300} problems (100 train, 200 test), confirming the README's ``305'' print statement is wrong \\
\code{dataset.split(...)} & evolution 65, holdout 200, rest 35: matches config exactly \\
\code{visualize\_test\_results.py} & Regenerates all 5 PNGs; prints tournament 0.7857, roulette 0.7717, elitism 0.7713, baseline 0.7900, matching the README table exactly \\
Recomputed means from the 4 JSONL logs & Independently reproduces the same four numbers; all logs have 200 entries, 0 infra failures \\
Recomputed win/loss from \code{comparison.json} & 26 wins, 22 losses, 152 ties: matches the README \\
Cost totals from \code{Costs/cost\_run\_*.json} & 13{,}002 calls, \$17.05: matches ``roughly 13,000'' and ``about \$17'' \\
\code{evaluator.run\_tests\_verbose} without \code{g++} & Returns \code{score = 0.0}, \code{reason = compile\_fail}: \textbf{contradicts} the README's ``rather than silently scoring zero'' \\
Best fitness monotonicity per \code{history.json} & tournament and roulette \textbf{non-monotonic and end below gen 0}; elitism monotonic \\
Winner present in any checkpoint & tournament \textbf{no}, roulette \textbf{no}, elitism yes \\
\code{grep} for genealogy in \code{*.py} & \textbf{No match}: the script is not committed \\
\code{grep} for chain-of-thought / random search & \textbf{No match}: neither baseline exists \\
\code{grep} for \code{TARGET\_TEMPERATURE} usage & Defined in \code{config.py}, \textbf{read nowhere} \\
README license claim vs LICENSE file & README says MIT; file is PolyForm Strict 1.0.0 \\
\midrule
\textbf{Not verified} & The full evolution loop, every Vertex AI call, and the compile-and-run fitness path. These need a billed GCP project and a \code{g++} install. All fitness numbers in this document are read from committed artifacts, not reproduced. \\
\bottomrule
\end{tabular}
\caption{Verification record. The repository's own analysis pipeline reproduces its headline
numbers exactly from the committed logs, which is a genuine strength: the results are
auditable without spending a dollar.}
\end{table}

\subsection{Everything labelled inferred in this document}

Stated plainly, so nothing here is mistaken for a confirmed fact:

\begingroup\sloppy
\begin{enumerate}
  \item \textbf{That the tournament run used elitism 0} despite its \code{config.json}
  recording 2. Inferred from the mathematical impossibility of an elite coexisting with
  decreasing best fitness. The artifacts do not record the value actually used.
  \item \textbf{That the 700-token operator cap interacts with the 2048-token thinking
  budget} to cause truncation. Consistent with the settings and with all three winners being
  truncated, but unverifiable without API access.
  \item \textbf{That \code{seeds.py}'s ten unused prompts are leftovers from the paper's
  population-15 ``full-scale setup''}. Inferred from the coincidence of the number 15.
  \item \textbf{That \code{sys.set\_int\_max\_str\_digits} is needed for a specific APPS
  problem's large numeric data}. Nothing records why it is there.
  \item \textbf{That the roulette/elitism comparison logs were made by hand-editing
  \code{run\_baseline\_eval.py}}. Inferred: the logs exist, the committed script cannot
  produce them, and no other script does.
\end{enumerate}
\endgroup

\newpage
%=============================================================================
\section{Summary: the honest scorecard}
%=============================================================================

\subsection{What is genuinely good}

\begin{itemize}
  \item \textbf{The fitness function.} Compile-and-run against real tests is the right
  answer and the hard one. It is what makes the negative result trustworthy rather than
  arguable.
  \item \textbf{Failure-conditioned mutation.} Feeding real compile errors and wrong outputs
  into the mutation operator is a genuinely good idea: a directed mutation in a
  gradient-free method.
  \item \textbf{RNG state checkpointing.} Serialising the full Mersenne Twister state so
  \code{--resume} is decision-identical is unusually careful.
  \item \textbf{Process-group killing.} \code{os.setsid} plus \code{os.killpg} is the
  correct way to bound untrusted subprocesses, and most first attempts get it wrong.
  \item \textbf{Cost accounting, including thinking tokens.} The \$17 figure is real and
  auditable, and folding thinking tokens in at the output rate is the detail that makes it
  match the bill.
  \item \textbf{Auditable results.} The analysis pipeline reproduces the headline numbers
  exactly from committed logs, offline. I checked; it does.
  \item \textbf{A README that already admits the headline negative result} and several of
  its own bugs. That is rarer than it should be.
\end{itemize}

\subsection{What to fix first}

Ordered by how much damage each does if someone else finds it first:

\begin{enumerate}
  \item \textbf{The license contradiction} (README says MIT, file is PolyForm Strict). One
  line. It is the only defect here with legal meaning.
  \item \textbf{The ``evolution reliably improved fitness'' claim} in the README, and the
  paper's matching claim. The data shows two of three runs \emph{regressing} below their
  seeds. Section~\ref{sec:regression}.
  \item \textbf{The paper's abstract}, which claims comparisons against chain-of-thought and
  random search that do not exist in the codebase.
  \item \textbf{Checkpoint the final generation} (one line). It silently breaks resume and
  it is the root cause of the broken genealogy.
  \item \textbf{Validate operator output for completeness.} All three winning prompts are
  truncated mid-sentence. Check the API's \code{finish\_reason}.
  \item \textbf{Add \code{g++} to the pre-flight check}, and fix the README sentence that
  describes this failure mode backwards.
  \item \textbf{Record the elitism value actually used} in each run's \code{config.json},
  not the module constant.
  \item \textbf{Commit the genealogy script}, or delete its orphaned outputs and drop the
  claim.
  \item \textbf{Parameterise \code{run\_baseline\_eval.py} by scheme}, so the four-way
  comparison is reproducible without editing source.
  \item \textbf{Statistical testing.} A 0.0043 gap over 200 problems needs a paired test or
  a bootstrap before it means anything in either direction. The README already says this.
  \item \textbf{Delete the dead code}: \code{run\_fitness\_batch}, \code{run\_tests},
  \code{reserved\_train\_problems}, the unused \code{PRICING} import, the Claude pricing
  rows, the dead temperature constants.
\end{enumerate}

\begin{keyidea}{The line to take in an interview}
Lead with the method, not the result. The project built an objective, expensive, honest
measuring instrument for a question that is usually answered by vibes, and then reported
that the answer was no. The instrument works: I re-ran its analysis and it reproduces its
own numbers exactly.

Then get ahead of the two things a careful reader will find: the winning prompts are all
truncated fragments, and two of three runs ended below their starting seeds because nothing
protected the best individual. Both have clean, short explanations. Both are more
interesting than the headline. Neither is in the write-up yet, and fixing that is the next
job.
\end{keyidea}

\end{document}
